\Needspace{12\baselineskip}
\subsubsection{3.2.20 UC20 查看每日排班与到岗情况}\label{uc-v3-20}
\begin{BSRUsecase}{UC20 查看每日排班与到岗情况}
主执行者 & 门店经理；工作人员查询本人安排 \\
辅执行者 & 无独立外部执行者；数据库、SSE 与计时模块均为系统内部机制。 \\
涉众利益 & 经理需要当天岗位覆盖和实际在岗；人员需要确认班次；隐私要求不得任意改写他人出勤。 \\
目标与触发 & 经理进入首页查看今天，或在每日排班中选择日期和门店。 \\
前置条件 & 会话有效；日期真实；访问范围由角色确定。 \\
成功后置条件 & 只读返回人员、计划班次、实际上下班和派生状态；不因查看产生打卡或调整库存。 \\
最小保证 & 失败不提交部分业务变更；权限或版本冲突返回明确原因。业务事实以服务器记录为准。 \\
基本流程 & \begin{minipage}[t]{\linewidth}\raggedright 1 执行者提交工作日期和可选门店范围。\\ 2 系统取得授权人员目录、该日期班次及出勤。\\ 3 系统按人员关联班次和记录，生成未排班、未打卡、在岗或已下班状态。\\ 4 系统按计划起止与实际时间标识迟到或早退提示。\\ 5 系统返回时间轴、到岗人数和班次明细。\\ 6 其他角色打卡或班次保存后，空闲页面重新读取同一日期数据。\strut\end{minipage} \\
扩展流程 & \begin{minipage}[t]{\linewidth}\raggedright 1a 日期无效：拒绝，不按服务器默认日期静默替换。\\ 2a 店员或技师：仅返回本人每日安排；经理、经营和管理员按当前授权查看概览。\\ 3a 无班次但有打卡：显示实际出勤，不抹去事实。\\ 3b 有班次无打卡：显示未打卡，不自动认定缺勤。\strut\end{minipage} \\
字段与状态事件 & date、storeId、person、shift、attendance、status、late、early；汇总人数从同一结果集计算。 \\
\rowcolor{GoldPale}业务规则 & BR-032 至 BR-034；等于计划边界不标记迟到早退；工作日期按挪威口径。 \\
质量与验收 & 验证日期和门店筛选、无班次、有班无打卡、在岗、下班、边界时间及跨角色同步。 \\
设计约束 & GET /api/roster；每日结果为查询投影，不额外保存每日排班快照表；经理首页默认挪威今天。 \\
\end{BSRUsecase}
